\title{Group Sequence Policy Optimization}

\begin{abstract}
We present Group Sequence Policy Optimization (GSPO), a robust, efficient, and high-performing reinforcement learning method tailored for large language model training. GSPO diverges from earlier approaches that rely on token-level importance ratios by instead establishing the importance ratio using the entire sequence likelihood, and subsequently applying sequence-level clipping, reward assignment, and optimization steps. Through extensive evaluation, GSPO demonstrates enhanced training efficiency and overall performance relative to the GRPO algorithm, significantly improves stability for Mixture-of-Experts (MoE) RL training scenarios, and offers potential for streamlining RL system architecture. These advantages of GSPO have played a pivotal role in the notable advances observed in the most recent Qwen3 model series.
\end{abstract}

\section{Introduction}
\label{sec:intro}

Reinforcement learning (RL) has become central to advancing the performance of language models \citep{o1,dpsk-r1,qwq32b,qwen3}. Employing large-scale RL techniques allows these models to address complex challenges, including those in competition mathematics and programming, by enabling more extensive and profound reasoning.

Achieving successful scaling of RL, particularly with increased computational resources, necessitates the preservation of stable and resilient training behaviors. Nevertheless, leading RL methods such as GRPO \citep{grpo} confront substantial stability challenges in training large-scale language models, frequently encountering disastrous and unrecoverable collapse \citep{qwen3, minimax-m1}. Such instability significantly impedes ongoing RL-driven enhancement of language model capabilities.

This work reveals that GRPO's instability arises from a fundamental flaw in how importance sampling weights are applied and interpreted within its algorithmic structure. The improper use leads to elevated training noise with significant variance, which intensifies as responses become longer. This effect is exacerbated by the clipping step, ultimately causing the model to collapse.

To overcome these foundational issues, we introduce \textbf{Group Sequence Policy Optimization (GSPO)}, a novel reinforcement learning algorithm for training large language models. The central contribution of GSPO is a rigorously defined importance ratio formulated in terms of sequence likelihood \citep{click}, conforming to the original tenets of importance sampling. Moreover, GSPO calculates normalized rewards as the advantages across multiple responses per query, thereby preserving consistency between reward computation at the sequence level and policy optimization.

Our experiments reveal that GSPO notably outperforms GRPO in terms of stability, efficiency, and overall training outcomes. Importantly, GSPO natively addresses the stability issues associated with reinforcement learning for large Mixture-of-Experts (MoE) architectures, obviating the reliance on elaborate stabilization methods and suggesting a route toward streamlining RL pipelines. These advantages have directly supported the marked progress observed in the most recent Qwen3 models. We anticipate that GSPO will serve as a reliable and scalable basis for further progress in large-scale reinforcement learning applied to language models.

\section{Preliminaries}

\paragraph{Notation}
Throughout this work, we represent an autoregressive language model with parameters $\theta$ as a policy $\pi_\theta$.
The variable $x$ refers to a query, while the collection of queries is denoted by $\mathcal{D}$. For any response $y$ generated in reply to query $x$, its probability under the policy $\pi_\theta$ is written as $\pi_\theta (y | x)=\prod_{t=1}^{|y|} \pi_\theta (y_t | x, y_{<t} )$, where $|y|$ indicates the token count within $y$. Each pair $(x, y)$, consisting of a query and its associated response, can be evaluated by a verifier $r$, which assigns a reward $r(x, y) \in [ 0, 1 ]$.

\paragraph{Proximal Policy Optimization (PPO)} Using samples generated from the old policy $\pi_{\theta_\text{old}}$, PPO \citep{ppo} constrains the policy update within a proximal region of the old policy through the clipping mechanism. Specifically, PPO employs the following objective for policy optimization (we omit the KL regularization term hereinafter for brevity, as it is not the focus of this paper):

\begin{align}
\mathcal{J}_\text{PPO}(\theta) = \mathbb{E}_{ x \sim \mathcal{D},\, y \sim \pi_{\theta_\text{old}}( \cdot | x) }
\left[ \frac{1}{|y|} \sum_{t=1}^{|y|} 
\min \left( w_{t}(\theta) \widehat{A}_{t},  \, \mathrm{clip} \left( w_{t}(\theta), 1 - {\varepsilon}, 1 + {\varepsilon}\right) \widehat{A}_{t} \right)
\right],
\end{align}

Here, the token importance ratio for $y_t$ is given by $
w_{t}(\theta) = \frac{ \pi_{\theta} (y_{t} | x, y_{<t}) }{ \pi_{\theta_\text{old}} (y_{t} | x,y_{<t})} $. The advantage $\widehat{A}_{t}$ of $y_t$ is approximated using an additional value model, and $\varepsilon$ denotes the clip threshold for these ratios.

A significant obstacle encountered with PPO in real-world applications is its substantial dependence on the value model. This model typically matches the policy model in scale, thereby imposing notable demands on both computational resources and memory. Additionally, the performance of PPO is intricately tied to the precision of the value estimates it produces. Obtaining a dependable value model is intrinsically difficult, and adapting it to support longer outputs and more intricate tasks only intensifies this difficulty.

\paragraph{Group Relative Policy Optimization (GRPO)} GRPO \citep{grpo} bypasses the need for the value model by computing the relative advantage of each response within a group of responses to the same query. Specifically, GRPO optimizes the following objective:

\begin{align}
\label{equ:grpo}
\mathcal{J}_\text{GRPO}(\theta) = \mathbb{E}_{ x \sim \mathcal{D},\, \{y_i\}_{i=1}^G \sim \pi_{\theta_\text{old}}( \cdot | x) }
\left[ \frac{1}{G} \sum_{i=1}^{G} \frac{1}{|y_i|} \sum_{t=1}^{|y_i|} 
\min \left( w_{i,t}(\theta) \widehat{A}_{i,t},  \, \mathrm{clip} \left( w_{i,t}(\theta), 1 - {\varepsilon}, 1 + {\varepsilon}\right) \widehat{A}_{i,t} \right)
\right],
\end{align}

where $G$ is the number of generated responses to each query $x$ (i.e., the group size), and the importance ratio $w_{i,t}(\theta)$ and advantage $\widehat{A}_{i,t}$ of token $y_{i,t}$ are:

\begin{align}
    w_{i,t}(\theta)=\frac{ \pi_{\theta} (y_{i,t} | x, y_{i,<t}) }{ \pi_{\theta_\text{old}} (y_{i,t} | x,y_{i,<t})},\quad\ 
    \widehat{A}_{i,t} = \widehat{A}_{i} = \frac{r(x, y_i) - \mathrm{mean} \left( \{ r(x, y_i) \}_{i=1}^G \right) }{ \mathrm{std} \left( \{ r(x, y_i) \}_{i=1}^G \right) },
\end{align}

respectively, where all the tokens in $y_i$ share the same advantage as $\widehat{A}_{i}$.

\section{Motivation}
\label{sec:motivation}

As model sizes increase and architectures become more sparse—particularly in the case of Mixture-of-Experts—and as response lengths grow, utilizing large rollout batch sizes is crucial for fully exploiting hardware resources during RL training. To enhance the efficiency of sample usage, practitioners commonly divide these extensive rollout batches into several mini-batches for conducting gradient updates. This approach naturally leads to an off-policy regime because the responses $y$ are generated under a previous policy $\pi_{\theta_\text{old}}$ instead of the policy $\pi_\theta$ currently undergoing optimization. Consequently, methods such as PPO and GRPO employ clipping strategies to curtail the influence of samples that are too divergent from the latest policy, thus ensuring more stable gradient calculations.

While mechanisms like clipping aim to manage this off-policy discrepancy, we identify a more fundamental issue in GRPO: \textit{its objective is ill-posed}. This problem becomes particularly acute when training large models on long-response tasks, leading to catastrophic model collapse. The ill-posed nature of the GRPO objective stems from a misapplication of importance sampling weights. The principle of importance sampling is to estimate the expectation of a function $f$ under a target distribution $\pi_\text{tar}$ by re-weighting samples drawn from a behavior distribution $\pi_\text{beh}$:

\begin{align}
\mathbb{E}_{z \sim \pi_\text{tar}} \left[ f(z) \right] = \mathbb{E}_{z \sim \pi_\text{beh}} \left[ \frac{ \pi_\text{tar} (z) }{ \pi_\text{beh} (z)} \, f(z) \right].
\end{align}

A key requirement for this approach is to aggregate over a substantial number of samples ($N \gg 1$) drawn from the behavior policy $\pi_\text{beh}$, enabling the importance weight $\frac{ \pi_\text{tar} (z) }{ \pi_\text{beh} (z)}$ to properly adjust for disparities between the sampled and target distributions.

By comparison, GRPO employs the importance weight $\frac{ \pi_{\theta} (y_{i,t} | x, y_{i,<t}) }{ \pi_{\theta_\text{old}} (y_{i,t} | x,y_{i,<t})}$ at each individual token step $t$. This weighting relies on only a single token $y_{i,t}$ sampled from each next-token conditional $\pi_{\theta_\text{old}}( \cdot | x, y_{i, <t})$, so it cannot reliably address the distribution mismatch. Rather than providing correction, it injects substantial noise into the gradient estimates; this noise accumulates throughout extended sequences and becomes worse due to the presence of the clipping mechanism. Our experiments confirm that such effects can result in catastrophic model collapse, from which recovery proves unattainable. Even restarting from previous checkpoints, making thorough adjustments to hyperparameters (such as clipping thresholds), increasing sequence length, or revising RL prompts fails to restore the model once collapse sets in.

These empirical findings indicate a central deficiency inherent in GRPO's methodology. The inadequacy of token-wise importance weighting reveals a more general lesson: \textbf{the unit at which the optimization criterion operates must align with the unit receiving reward}. Rewards are assigned per sequence, so correcting for off-policy discrepancies at the token level is inherently flawed. This recognition leads us to abandon the token-level objective in favor of sequence-level optimization, with importance weights applied directly at the sequence granularity.

\section{Algorithm}

\subsection{GSPO: Group Sequence Policy Optimization}

Although the use of the token-level importance weight $\frac{ \pi_{\theta} (y_{i,t} | x, y_{i,<t}) }{ \pi_{\theta_\text{old}} (y_{i,t} | x,y_{i,<t})}$ poses issues in GRPO, we note that, for language generation tasks, the \textit{sequence-level} importance weight $\frac{ \pi_{\theta} (y | x) }{ \pi_{\theta_\text{old}} (y | x)}$ possesses a well-defined theoretical significance. Specifically, it quantifies the degree to which a response $y$, drawn from $\pi_{\theta_\text{old}} (\cdot | x)$, differs from the current policy $\pi_{\theta} (\cdot | x)$. This correspondence is in line with sequence-level rewards, and it provides a meaningful metric for gauging the efficacy of the clipping process.

Based on this straightforward observation, we propose the \textbf{Group Sequence Policy Optimization (GSPO)} algorithm. GSPO employs the following sequence-level optimization objective:

\begin{align}
\mathcal{J}_\text{GSPO} (\theta) =
\mathbb{E}_{ x \sim \mathcal{D},\, \{y_i\}_{i=1}^G \sim \pi_{\theta_\text{old}}( \cdot | x) }
\left[ 
\frac{1}{G} \sum_{i=1}^{G}
\min \left( s_{i}(\theta)  \widehat{A}_{i},  \, \mathrm{clip} \left( s_{i}(\theta), 1 - {\varepsilon}, 1 + {\varepsilon} \right) \widehat{A}_{i} \right) 
\right],
\label{equ:gspo}
\end{align}

where we adopt the group-based advantage estimation:

\begin{align}
\widehat{A}_{i} = \frac{r(x, y_i) - \mathrm{mean} \left( \{ r(x, y_i) \}_{i=1}^G \right) }{ \mathrm{std} \left( \{ r(x, y_i) \}_{i=1}^G \right) },
\end{align}

and define the importance ratio $s_{i}(\theta)$ based on sequence likelihood \citep{click}:

\begin{align}
s_{i}(\theta) = \left( \frac{ \pi_{\theta} (y_i | x) }{ \pi_{\theta_\text{old}} (y_i | x)} \right)^{\frac{1}{|y_i|}}
=
\exp \left( \frac{1}{|y_i|} \sum_{t=1}^{|y_i|} \log \frac{ \pi_{\theta} (y_{i,t} | x, y_{i,<t}) }{ \pi_{\theta_\text{old}} (y_{i,t} | x,y_{i,<t})} \right).
\end{align}

Consequently, GSPO implements clipping at the response level rather than on single tokens, thereby removing excessively ``off-policy'' samples from the gradient computation. This approach aligns with both reward assignment and optimization processes at the sequence level. Additionally, we utilize length normalization within $s_{i}(\theta)$ to mitigate variance and ensure $s_{i}(\theta)$ remains within a consistent numerical interval. In the absence of such normalization, changes in token likelihoods could substantially impact the importance ratio calculated at the sequence level, and responses of varying lengths would necessitate distinct clipping intervals. Moreover, it is worth noting that the clipping thresholds employed by GSPO diverge in scale from those used by earlier methods such as GRPO, reflecting fundamental differences in the way importance ratios are defined.

\subsection{Gradient Analysis}
\label{subsec:gradient}

We can derive the gradient of the GSPO objective as follows (clipping is omitted for brevity):

\begin{align}
\nabla_{\theta} \mathcal{J}_\text{GSPO} (\theta)
=&\ 
\nabla_{\theta} \mathbb{E}_{ x \sim \mathcal{D},\, \{y_i\}_{i=1}^G \sim \pi_{\theta_\text{old}}( \cdot | x) }
\left[ 
\frac{1}{G} \sum_{i=1}^{G} s_{i}(\theta) \widehat{A}_{i}
\right] \\
= &\ 
\mathbb{E}_{ x \sim \mathcal{D},\, \{y_i\}_{i=1}^G \sim \pi_{\theta_\text{old}}( \cdot | x) }
\left[ 
\frac{1}{G} \sum_{i=1}^{G}
s_{i}(\theta) \widehat{A}_{i}
\cdot \nabla_{\theta} \log s_{i}(\theta)
\right] \\
=&\ 
\mathbb{E}_{ x \sim \mathcal{D},\, \{y_i\}_{i=1}^G \sim \pi_{\theta_\text{old}}( \cdot | x) }
\left[ 
\frac{1}{G} \sum_{i=1}^{G} 
\left( \frac{ \pi_{\theta} (y_i | x) }{ \pi_{\theta_\text{old}} (y_i | x)} \right)^{\frac{1}{|y_i|}} \widehat{A}_{i} 
\cdot \frac{1}{|y_i|} \sum_{t=1}^{|y_i|} \nabla_{\theta} \log \pi_{\theta} (y_{i,t} | x, y_{i,<t}) 
\right].
\label{equ:gspo_grad}
\end{align}

For comparison, the gradient of the GRPO objective is as follows (note that $\widehat{A}_{i,t} = \widehat{A}_{i}$):

\begin{align}
\nabla_{\theta} \mathcal{J}_\text{GRPO} (\theta) 
=&\ 
\nabla_{\theta} \mathbb{E}_{ x \sim \mathcal{D},\, \{y_i\}_{i=1}^G \sim \pi_{\theta_\text{old}}( \cdot | x) }
\left[ \frac{1}{G} \sum_{i=1}^{G} \frac{1}{|y_i|} \sum_{t=1}^{|y_i|} 
w_{i,t}(\theta) \widehat{A}_{i,t}
\right] \\
=&\
\mathbb{E}_{ x \sim \mathcal{D},\, \{y_i\}_{i=1}^G \sim \pi_{\theta_\text{old}}( \cdot | x) }
\left[ \frac{1}{G} \sum_{i=1}^{G} \widehat{A}_{i} 
\cdot \frac{1}{|y_i|} \sum_{t=1}^{|y_i|} 
\frac{ \pi_{\theta} (y_{i,t} | x, y_{i,<t}) }{ \pi_{\theta_\text{old}} (y_{i,t} | x,y_{i,<t})} 
\nabla_{\theta} \log \pi_{\theta} (y_{i,t} | x, y_{i,<t})  
\right].
\end{align}

As a result, the primary difference separating GSPO and GRPO is centered on \textit{the manner in which each method assigns weights to the gradients of token log likelihoods}. Specifically, GRPO applies an ``importance weight'' to tokens, calculated as $\frac{ \pi_{\theta} (y_{i,t} | x, y_{i,<t}) }{ \pi_{\theta_\text{old}} (y_{i,t} | x,y_{i,<t})}$. These variable weights can span the ranges $(0, 1+\varepsilon]$ for cases where $\widehat{A}_{i} > 0$, or $[1-\varepsilon, +\infty)$ when $\widehat{A}_{i} < 0$, and their non-uniformity is significant. As training continues, such disparities may accumulate and induce unpredictable effects. By contrast, GSPO employs uniform weighting for all tokens within a response, thereby removing the source of instability that arises in GRPO.

\subsection{GSPO-token: A Token-level Objective Variant}

In scenarios like multi-turn RL, we may desire a finer-grained advantage adjustment than the sequence level. To this end, we introduce a token-level objective variant of GSPO, namely \textbf{GSPO-token}, to allow token-wise advantage customization:

\begin{align}
\mathcal{J}_\text{GSPO-token}(\theta)
=
\mathbb{E}_{ x \sim \mathcal{D},\, \{y_i\}_{i=1}^G \sim \pi_{\theta_\text{old}}( \cdot | x) }
\left[ 
\frac{1}{G} \sum_{i=1}^{G} \frac{1}{|y_i|} \sum_{t=1}^{|y_i|} 
\min \left( s_{i,t}(\theta) \widehat{A}_{i,t},  \, \mathrm{clip} \left( s_{i,t}(\theta), 1 - {\varepsilon}, 1 + {\varepsilon} \right) \widehat{A}_{i,t} \right) 
\right],
\label{equ:gspo-token}
\end{align}

where

\begin{align}
s_{i,t}(\theta) 
= \mathrm{sg} \left[ s_{i}(\theta) \right]  \cdot \frac{ \pi_{\theta} (y_{i,t} | x, y_{i,<t}) }{ \mathrm{sg} \left[ \pi_{\theta} (y_{i,t} | x, y_{i,<t}) \right] },
\end{align}

and $\mathrm{sg}[\cdot]$ denotes only taking the numerical value but stopping the gradient, corresponding to the \texttt{detach} operation in PyTorch. The gradient of GSPO-token can be derived as:

\begin{align}
\nabla_{\theta} \mathcal{J}_\text{GSPO-token}(\theta)
=&\ 
\nabla_{\theta} \mathbb{E}_{ x \sim \mathcal{D},\, \{y_i\}_{i=1}^G \sim \pi_{\theta_\text{old}}( \cdot | x) }
\left[ 
\frac{1}{G} \sum_{i=1}^{G}
\frac{1}{|y_i|} \sum_{t=1}^{|y_i|} 
s_{i,t}(\theta) \widehat{A}_{i,t}
\right] \\
=&\ 
\mathbb{E}_{ x \sim \mathcal{D},\, \{y_i\}_{i=1}^G \sim \pi_{\theta_\text{old}}( \cdot | x) }
\left[ 
\frac{1}{G} \sum_{i=1}^{G} s_i (\theta) \cdot \frac{1}{|y_i|} \sum_{t=1}^{|y_i|} 
\widehat{A}_{i,t} \frac{ \nabla_{\theta} \pi_{\theta} (y_{i,t} | x, y_{i,<t}) }{ \pi_{\theta} (y_{i,t} | x, y_{i,<t}) }
\right] \\
=&\ 
\mathbb{E}_{ x \sim \mathcal{D},\, \{y_i\}_{i=1}^G \sim \pi_{\theta_\text{old}}( \cdot | x) }
\left[ 
\frac{1}{G} \sum_{i=1}^{G}  \left( \frac{ \pi_{\theta} (y_i | x) }{ \pi_{\theta_\text{old}} (y_i | x)} \right)^{\frac{1}{|y_i|}}
\cdot \frac{1}{|y_i|} \sum_{t=1}^{|y_i|}
\widehat{A}_{i,t} \nabla_{\theta} \log \pi_{\theta} (y_{i,t} | x, y_{i,<t})  
\right].
\label{equ:gspo-token_grad}
\end{align}

Observe that the expression $\frac{ \pi_{\theta} (y_{i,t} | x, y_{i,<t}) }{ \mathrm{sg} \left[ \pi_{\theta} (y_{i,t} | x, y_{i,<t}) \right] }$ evaluates to 1, implying that $s_{i,t}(\theta)$ is, in practice, equivalent to $s_{i}(\theta)$ in terms of numerical output. By examining Equation~\eqref{equ:gspo} alongside Equation~\eqref{equ:gspo-token}, and Equation~\eqref{equ:gspo_grad} with Equation~\eqref{equ:gspo-token_grad}, it is evident that GSPO-token and GSPO coincide exactly in their optimization objective, clipping criterion, and analytical gradient, as long as the advantage values assigned to every token in the output $y_i$ are made uniform (specifically, $\widehat{A}_{i,t} = \widehat{A}_{i}$ for all $t$). However, GSPO-token provides greater adaptability, permitting advantage assignment to be customized on a per-token basis.

\section{Experiments and Discussion}

\subsection{Empirical Results}

We conduct experiments using a cold-start model, fine-tuned from Qwen3-30B-A3B-Base, and present both the training reward trajectories and evaluation metrics on the AIME'24 benchmark (measured by average Pass@1 with 32 samples), LiveCodeBench (202410-202502, average Pass@1 with 8 samples), and CodeForces (using Elo Rating). Throughout reinforcement learning training, the rollout data for each batch is subdivided into four mini-batches to perform gradient updates. For GSPO, the left and right clipping boundaries in Equation~\eqref{equ:gspo} are chosen as 3e-4 and 4e-4, accordingly. As a point of comparison, we use GRPO as the baseline and configure its left and right clipping thresholds in Equation~\eqref{equ:grpo} to 0.2 and 0.27, respectively—values determined through careful tuning to maintain experimental fairness. Importantly, GRPO requires application of the Routing Replay training method to achieve stable convergence in MoE RL, which we will further elaborate on in \S~\ref{subsec:routing_replay}, whereas \textit{GSPO does not require such a strategy}.

As illustrated in Figure~\ref{fig:results}, GSPO-based training exhibits consistent stability across all stages. Our analysis indicates that \textbf{GSPO enables ongoing enhancements in performance by scaling up training compute, systematically refreshing the query set, and lengthening the generated output}. Additionally, GSPO demonstrates greater training efficiency relative to GRPO, achieving superior accuracy and benchmark scores with equivalent compute resources and query consumption. Lastly, we have implemented GSPO in the reinforcement learning training of the most recent Qwen3 models, providing compelling evidence that GSPO effectively harnesses RL scaling to improve large language model capabilities.

\subsection{Curious Observation on Clipping Fractions}

GSPO diverges from GRPO primarily in its method of clipping, opting to truncate entire response sequences as opposed to individual tokens. As illustrated in Figure~\ref{fig:clipping}, there is a dramatic difference—amounting to two orders of magnitude—in the proportion of tokens subjected to clipping between GSPO and GRPO; moreover, adjusting the clipping thresholds does not eliminate this pronounced discrepancy. Nonetheless, GSPO manages to attain greater training efficiency than GRPO even though it clips a substantially higher number of tokens, thereby incorporating fewer tokens during the training process (or when estimating gradients). This somewhat surprising result—that more extensive token clipping enhances training efficiency—implies that GRPO’s approach of estimating gradients at the token level introduces considerable noise and is suboptimal for extracting informative samples. GSPO, by adopting sequence-level clipping, yields a more dependable and robust learning signal.

\subsection{Benefit of GSPO for MoE Training}
\label{subsec:routing_replay}

\paragraph{Background}
Unlike the RL training of conventional dense models, the inherently sparse activation patterns in MoE architectures give rise to distinct issues with training stability. Notably, our investigations indicate that using the GRPO algorithm with MoE models brings about significant \textit{expert-activation volatility}, which can obstruct successful RL convergence. Concretely, we observe that after applying one or more rounds of gradient updates, the specific set of experts activated for a fixed response may vary considerably. For instance, experiments with the 48-layer Qwen3-30B-A3B-Base model show that, following each RL gradient step and given the same sampled rollout, approximately 10\% of the experts chosen by the updated policy $\pi_{\theta}$ do not overlap with those activated by the previous policy $\pi_{\theta_\text{old}}$.
This effect, which is magnified in larger and deeper MoE models, leads to substantial swings in the token-level importance ratios $w_{i,t}(\theta) = \frac{ \pi_{\theta} (y_{i,t} | x, y_{i,<t}) }{ \pi_{\theta_\text{old}} (y_{i,t} | x,y_{i,<t})}$. Such volatility degrades their reliability, as previously examined in \S~\ref{sec:motivation} and~\ref{subsec:gradient}, thereby impeding the RL training process from achieving proper convergence.

\paragraph{Our Previous Approach} To address this issue, our earlier method involved the \textbf{Routing Replay} training paradigm. Concretely, we stored the experts selected by $\pi_{\theta_\text{old}}$ and enforced the same routing decisions when evaluating $\pi_{\theta}$ during the calculation of importance weights, $w_{i,t}(\theta) = \frac{ \pi_{\theta} (y_{i,t} | x, y_{i,<t}) }{ \pi_{\theta_\text{old}} (y_{i,t} | x,y_{i,<t})}$. This guarantees that, for each token $y_{i,t}$, both $\pi_{\theta} (y_{i,t} | x, y_{i,<t})$ and $\pi_{\theta_\text{old}} (y_{i,t} | x,y_{i,<t})$ utilize identical expert activations. As a result, this approach stabilizes the computation of token-level importance ratios, enabling consistent optimization with respect to the same active network throughout the course of gradient-based updates. 
As illustrated in Figure~\ref{fig:routing_replay}, Routing Replay is critical for achieving proper convergence during GRPO training of MoE architectures.

\paragraph{Benefit of GSPO} While Routing Replay ensures that GRPO training of MoE models converges as intended, its method of reusing routing configurations introduces extra demands on memory and communication, and may restrict the effective utilization of the MoE model’s potential. On the other hand, as illustrated in Figure~\ref{fig:results}, GSPO avoids reliance on Routing Replay. GSPO straightforwardly computes the importance ratios $s_i(\theta)$ in the typical manner, achieving stable optimization and expected convergence. The primary observation is that GSPO evaluates only the sequence likelihood (i.e., $\pi_{\theta} (y_i | x)$) and does not depend on the specific token likelihoods (i.e., $\pi_{\theta} (y_{i,t} | x, y_{i,<t})$), meaning GSPO is robust to fluctuations at the token level. Because the MoE architecture consistently preserves its ability to model language, sequence likelihood remains relatively stable. In essence, GSPO directly addresses the problem of expert-activation instability within MoE models, removing the necessity for elaborate techniques such as Routing Replay. This both reduces complexity and enhances stability during training, empowering the model to make full use of its inherent capacity without imposed limitations.

\subsection{Benefit of GSPO for RL Infrastructure}

Due to variations in numerical precision between training platforms (such as Megatron) and inference platforms (such as SGLang and vLLM), it is common practice to rely on the training engine to recalculate the likelihoods of sampled responses under the prior policy $\pi_{\theta_\text{old}}$. Nevertheless, GSPO bases its optimization on sequence-level likelihoods instead of the more sensitive token-level ones. This approach inherently provides greater robustness against precision inconsistencies. As a result, GSPO enables the direct utilization of likelihoods provided by the inference engine for optimization purposes, eliminating the overhead associated with re-evaluating them using the training engine. This capability proves particularly advantageous for use cases like partial rollout, multi-turn RL, and training-inference separated systems.

\section{Conclusion}

Group Sequence Policy Optimization (GSPO) is introduced as a novel reinforcement learning method tailored for the training of large language models. GSPO leverages the core idea of importance sampling, using importance ratios determined by sequence likelihood, and applies sequence-level techniques such as clipping, rewarding, and optimization. Empirical results show GSPO offers markedly better training stability, efficiency, and overall performance relative to GRPO, particularly excelling in large-scale reinforcement learning settings involving MoE models. These advancements serve as the basis for the significant improvements found in recent Qwen3 models. By establishing GSPO as a robust and scalable foundation, ongoing efforts aim to further scale reinforcement learning and drive substantial progress in artificial intelligence.

\end{document}